\documentclass[12pt]{article}
\usepackage[a4paper, margin=1in]{geometry}
\usepackage{graphicx}
\usepackage{amsmath}
\usepackage{amssymb}
\usepackage{dsfont}
\usepackage{comment}
\usepackage{xcolor}

\newcommand{\outdated}[1]{{\color{gray} #1}}
\newcommand{\todo}[1]{{\color{red} #1}}

\begin{document}

\section{Notation}
We model four events each year: growth, colonization, fizzle, and end-of-wave impact.

Let $(k,t)$ be the time index, where $k$ represents the year index and $t$ represents the within-year event index:
$t=0$ for the time at the beginning of the year,
$t=1$ for the time after growth,
$t=2$ for the time after colonization,
$t=3$ for the time after fizzle,
and $t=4$ for the time at the end of the year (after the wave).
We interpret $(k,0)$ as the same time as $(k-1,4)$.

Assume there are $n$ patches and we label them with $p_i$ ($i=1,\dots,n$).
For each patch we use $(N_i^{(k,t)}, I_i^{(k,t)})$ to describe its state.
$N_i^{(k,t)}$ represents population size in patch $p_i$ at time $(k,t)$, and
\[
I_i^{(k,t)}\in\{0,1,2,\dots\}
\]
is the number of infected hosts in patch $p_i$ at time $(k,t)$.
We define the (derived) infection indicator
\[
\mathds{1}_i^{(k,t)}=\mathds{1}\{I_i^{(k,t)}\ge 1\}.
\]

Let $U \sim \text{Uniform}(0, 1)$ be a random variable uniformly distributed over $(0,1)$, and we define
\[
\mathds{1}_{U \leq P} =
\begin{cases}
1, & \text{if } U \leq P, \\
0, & \text{otherwise.}
\end{cases}
\]

\section{Parameters}
Host:
\begin{itemize}
\item $r$: growth rate
\item $K$: carrying capacity of a patch (all patches are of equal area)
\item $c\in[0,1]$: probability that an individual host moves to another patch during colonization
\end{itemize}

Epidemic:
\begin{itemize}
\item $\mathcal{R}_0$: basic reproduction number
\item $P_f = 1/\mathcal{R}_0$: fizzle probability for a \emph{single} introduced infection
\item $z=z(\mathcal{R}_0)$: final size fraction of a non-fizzling within-patch epidemic wave
\item $D$: disease-induced mortality fraction among infected hosts (assume $D=1$)
\end{itemize}

Propagule (infectious fleas):
\begin{itemize}
\item $\nu$: scaling constant linking the previous year's epidemic size in a patch to the total number of infectious propagules generated
\item $\rho$: mean number of infectious propagules that a single migrating live host can carry (carrying capacity per host)
\end{itemize}

Infection:
\begin{itemize}
\item $\beta \ge 0$: mean number of effective bites per propagule during the colonization season
\item $p\in[0,1]$: probability that a single effective bite infects the bitten host 
\item $\alpha=\beta p$: transmission efficiency (successful infections) per propagule
\end{itemize}

\paragraph{Propagule assumption}
We assume free fleas do not survive beyond the current season, so all propagules relevant for year $k$ are generated by the year $k-1$ wave.

\section{Growth}
We assume that reproduction of rats follows logistic growth:
\begin{equation}
N_i^{(k,1)}=N_i^{(k,0)}+rN_i^{(k,0)}\left(1-\frac{N_i^{(k,0)}}{K}\right).
\end{equation}
Assume infection counts do not change during growth:
\begin{equation}
I_i^{(k,1)} = I_i^{(k,0)}=0.
\end{equation}

\section{Colonization and Fizzle}
\subsection{Population size}
We assume that each individual rat moves with probability $c$ and lands in a random patch.
We ignore distances and assume that all patches are equally available to a moving rat:
\begin{equation}
N_i^{(k,2)}=(1-c)N_i^{(k,1)}+\frac{c\sum_{j=1}^n N_j^{(k,1)}}{n}.
\end{equation}
The first term is the population staying in the patch, and the second term is the population moving into the patch.

\subsection{Propagules generated by the previous year}
Let $\mathds{1}_j^{(k-1,3)}=\mathds{1}\{I_j^{(k-1,3)}\ge 1\}$ indicate whether patch $p_j$ is infected after the fizzle step in year $k-1$.
The number of dead and infected hosts produced by the wave in patch $j$ in year $k-1$ is 
\begin{equation}
S_j^{(k-1)} = zD\,N_j^{(k-1,3)}\,\mathds{1}_j^{(k-1,3)}.
\label{eq:Sjk}
\end{equation}

\paragraph{Propagule capacity and distribution}
The total expected number of infectious propagules generated in patch $j$ is driven entirely by dead hosts, which is $\nu S_j^{(k-1)}$. 
However, cross-patch transmission is strictly driven by migrating live hosts. The number of live hosts migrating from patch $j$ is $c N_j^{(k,1)}$. These migrating hosts can carry a maximum expected number of $\rho c N_j^{(k,1)}$ propagules.

\outdated{test}
We define the expected number of exported propagules $\lambda_{E,j}^{(k)}$ and retained propagules $\lambda_{R,j}^{(k)}$ as:
\begin{align}
\lambda_{E,j}^{(k)} & = \phi \left( \nu S_j^{(k-1)}, \, \rho c N_j^{(k,1)} \right)   \label{eq:lambdaE} \\
\lambda_{R,j}^{(k)} &= \nu S_j^{(k-1)} - \lambda_{E,j}^{(k)} \label{eq:lambdaR}
\end{align}
with $\phi(F,R) \equiv F \left( 1- \exp(-R/R) \right)$ (a ``min-like'' function that approaches $F$ smoothly from below as $R \to \infty$ and \emph{vice versa})
\outdated{we previously defined $\lambda_{E,j}^{(k)}$ as $\min\left( \nu S_j^{(k-1)}, \, \rho c N_j^{(k,1)} \right)$}

By Poisson thinning of the total propagule pool, the actual number of retained propagules $R_j^{(k)}$ and total exported propagules $E_j^{(k)}$ generated by patch $j$ are:
\begin{align}
R_j^{(k)} &\sim \mathrm{Poisson}\!\left(\lambda_{R,j}^{(k)}\right) \label{eq:Rjk} \\
E_j^{(k)} &\sim \mathrm{Poisson}\!\left(\lambda_{E,j}^{(k)}\right) \label{eq:Ejk}
\end{align}

Exported propagules are assumed to land uniformly at random among the $n$ patches.
The number of propagules exported from source patch $j$ that arrive in target patch $i$ is
\begin{equation}
Y_{j\to i}^{(k)} \sim \mathrm{Poisson}\!\left(\frac{\lambda_{E,j}^{(k)}}{n}\right).
\label{eq:Yjtoi}
\end{equation}
Summing over all sources gives the total number of imported propagules in patch $i$:
\begin{equation}
Y_{i,\mathrm{import}}^{(k)} = \sum_{j=1}^n Y_{j\to i}^{(k)}.
\label{eq:Yimport}
\end{equation}

\paragraph{Total propagules present in patch $i$}
The total number of propagules present in patch $i$ at colonization is
\begin{equation}
G_i^{(k)} = R_i^{(k)} + Y_{i,\mathrm{import}}^{(k)}.
\label{eq:Gi}
\end{equation}

\subsection{Newly infected hosts}
We now derive the per-host infection probability and the distribution of the number of newly infected hosts created during colonization.

\paragraph{Per-host infection probability.}
Fix a particular host in patch $i$. Given $G_i^{(k)}$ propagules, assume:
(i) each propagule makes a Poisson($\beta$) number of effective bites;
(ii) each bite chooses a host uniformly among the $N_i^{(k,2)}$ hosts in the patch;
(iii) each bite infects with probability $p$, independently.

Then the number of \emph{successful infection attempts} on the focal host is Poisson with mean
\[
\frac{\beta p\,G_i^{(k)}}{N_i^{(k,2)}} = \frac{\alpha\,G_i^{(k)}}{N_i^{(k,2)}}.
\]
Therefore, the probability that the focal host becomes infected is
\begin{equation}
\pi_i^{(k)} = 1-\exp\!\left(-\frac{\alpha\,G_i^{(k)}}{N_i^{(k,2)}}\right)\approx\frac{\alpha\,G_i^{(k)}}{N_i^{(k,2)}}.
\label{eq:pi}
\end{equation}

\paragraph{Number of newly infected hosts}
Given $\pi_i^{(k)}$, the number of hosts infected is
\[
X_i^{(k)} \mid \left(G_i^{(k)},N_i^{(k,2)}\right) \sim \mathrm{Binomial}\!\left(N_i^{(k,2)},\,\pi_i^{(k)}\right).
\]
When $\pi_i^{(k)}\ll 1$, we can approximate this using the expectation of the total propagules present in the patch:

\begin{equation}
\begin{aligned}
X_i^{(k)} 
&\approx \mathrm{Poisson}\!\left(N_i^{(k,2)}\,\pi_i^{(k)}\right) \\
&= \mathrm{Poisson}\!\left(\alpha\,G_i^{(k)}\right) \\
&\approx \mathrm{Poisson}\!\left( \alpha \left[ \lambda_{R,i}^{(k)} + \frac{1}{n}\sum_{j=1}^n \lambda_{E,j}^{(k)} \right] \right)
\end{aligned}
\label{eq:X_poisson_approx}
\end{equation}

\paragraph{Update infected count after colonization.}
We define
\begin{equation}
I_i^{(k,2)} = I_i^{(k,1)} + X_i^{(k)} = X_i^{(k)}.
\end{equation}

\section{Fizzle}

Let $P_f=1/\mathcal{R}_0$ be the probability that a single introduction fails to establish.
If $X$ introductions occur independently, then the probability that all of them fizzle is $P_f^X$. Therefore,
\[
P(\text{establish}) = 1-\mathbb{E}[P_f^X].
\]

Since $X\sim \mathrm{Poisson}(\lambda)$, 
\[
\mathbb{E}[P_f^X]
= \sum_{x=0}^{\infty} P_f^x \frac{e^{-\lambda}\lambda^x}{x!}
= e^{-\lambda} \sum_{x=0}^{\infty}\frac{(\lambda P_f)^x}{x!}
= e^{-\lambda} e^{\lambda P_f}
= e^{-\lambda(1-P_f)}.
\]
Therefore
\[
P(\text{establish}) = 1-e^{-\lambda(1-P_f)}
\]
where the expected number of successful primary infections $\lambda$ is:
\[
\lambda = \alpha \left[ \lambda_{R,i}^{(k)} + \frac{1}{n}\sum_{j=1}^n \lambda_{E,j}^{(k)} \right]
\]

Population size does not change:
\begin{equation}
N_i^{(k,3)} = N_i^{(k,2)}.
\end{equation}

\section{Major epidemic}
In patches that experience a major epidemic, the population is reduced by the epidemic wave:
\begin{equation}
N_i^{(k,4)}=\left(1-zD\,\mathds{1}_i^{(k,3)}\right)N_i^{(k,3)}.
\end{equation}

We assume infected rats die quickly and no infected rats remain at the end of a wave:
\begin{equation}
I_i^{(k,4)} = 0 \qquad \text{for all } i.
\end{equation}

\section{Year-to-year update}
\begin{equation}
N_i^{(k+1,0)} = N_i^{(k,4)},\qquad
I_i^{(k+1,0)} = I_i^{(k,4)}=0.
\end{equation}

\end{document}
